% Laget automatisk av elevark/lag_elevark.py fra sim/to-baner.html. Endringer her blir overskrevet.
\documentclass[11pt]{article}
\usepackage{elevark}
\begin{document}
\elevarkhode{Fysikk 1 · Mekanisk energi}{To baner med samme fallhøyde}{https://lektorbo.no/sim/to-baner.html}{Simuleringen}

To like vogner starter i ro i samme høyde og glir ned hver sin bane til samme mål. Bane A er rett. Bane B er krum.

\section*{Oppgaver}

Jobb gjerne i par. Skriv ned gjetningen og begrunnelsen \emph{før} dere trykker Start.

\begin{oppgaveliste}
\oppgave \fase{gjett}{Gjett}Bruk $h = 1{,}00\,\mathrm{m}$, $μ = 0$ og «Bratt start» på bane B. Hvilken vogn kommer først i mål? Hvilken vogn har størst fart i mål?
\svar
\oppgave \fase{test}{Test}Kjør simuleringen. Stemte tipsene? Se på fartsgrafen: Hvilken vogn har størst fart underveis?
\svar
\oppgave \fase{forklar}{Forklar}Bruk energibevaring til å regne ut farten i mål, og sammenlign med tabellen. Forklar hvorfor B kommer først selv om farten i mål er den samme.
\svar
\oppgave \fase{gjett}{Gjett}Hva skjer med tiden, farten og energien hvis massen økes til 5,0~kg? Test og forklar med likningen fra oppgave 3.
\svar
\oppgave \fase{gjett}{Gjett}Bytt bane B til «Slak start». Hvem vinner nå? Test, og forklar ved hjelp av fartsgrafen.
\svar
\oppgave \fase{forklar}{Forklar}Sett $μ = 0{,}10$. Vis at friksjonsarbeidet på den rette banen er $W_{\mathrm{R}} = μmgL$, der $L = 2{,}0\,\mathrm{m}$ er den horisontale lengden. Regn ut farten i mål for vogn A, og kontroller med simuleringen.
\svar
\oppgave \fase{utfordring}{Utfordring}Med friksjon blir friksjonsarbeidet ulikt på de to banene, selv om begge er 2,0~m horisontalt. Hvorfor? Tips: Hvor stor er normalkraften i bunnen av en bue?
\svar
\end{oppgaveliste}

\end{document}
